%%%%%%%%%%%%%%%%%%%%%%%%%%%%%%%%%
%
%       EXERCÍCIO
%
%%%%%%%%%%%%%%%%%%%%%%%%%%%%%%%%%

\ifdefstring{\atividade}{prova}{%
    \renewcommand{\valorquestao}{\ValorQ\ pontos}
}%

\begin{exercicioBanco}[\valorquestao]
    Sorteia-se ao acaso um dentre os números 9, 12, 18 e 27 e é feita a decomposição do número sorteado em fatores primos. Sejam \(D\) e \(T\), as variáveis que representam, respectivamente, o número de vezes em que o 2 e o 3 aparecem na decomposição. Obtenha a tabela de frequência conjunta de \(D\) e \(T\).
\end{exercicioBanco}

